\QHeading{第1题　能围成三角形的在（　）里画“$\checkmark$”。}
\par\noindent


\PlainItem{① 3 cm　4 cm　6 cm（{\color{answerred}$\checkmark$}）}
\par

\PlainItem{② 1 cm　2 cm　3 cm（{\color{answerred}$\times$}）}
\PlainItem{③ 5 cm　7 cm　11 cm（{\color{answerred}$\checkmark$}）}
\PlainItem{④ 12 m　15 m　29 m（{\color{answerred}$\times$}）}
\par\noindent

\QHeading{第2题　下面每组中的3根小棒能围成三角形吗? 能的在（　）里画“$\checkmark$”。}

\par\noindent\makebox[\linewidth][c]{
\begin{tikzpicture}[x=1bp,y=-1bp,color=black,draw=black,text=black]
\begin{scope}\clip (56.0,242.0) rectangle (543.0,374.0);
\hdhead{42}{多边形的认识——信息窗2}
\hdsection{1}{140}{能围成三角形的在（　）里画“$\checkmark$”。}
\q{103}{197}{270}{12.5}{① 3 cm　4 cm　6 cm}
\q{480}{197}{270}{12.5}{② 1 cm　2 cm　3 cm}
\q{103}{251}{270}{12.5}{③ 5 cm　7 cm　11 cm}
\q{480}{251}{270}{12.5}{④ 12 m　15 m　29 m}
\q{355}{195}{28}{13}{（}
\q{409}{195}{25}{13}{）}
\q{736}{195}{28}{13}{（}
\q{790}{195}{25}{13}{）}
\q{355}{250}{28}{13}{（}
\q{409}{250}{25}{13}{）}
\q{736}{250}{28}{13}{（}
\q{790}{250}{25}{13}{）}
\hdsection{2}{314}{下面每组中的3根小棒能围成三角形吗? 能的在（　）里画“$\checkmark$”。}
\begin{scope}[scale=.6666667,draw=black,text=black]\path[draw=black,line width=.45pt,fill=bookpink] (102,390) rectangle (210,398);\draw[white,line width=.7pt] (104,392)--(208,392);\end{scope}
\begin{scope}[scale=.6666667,draw=black,text=black]\node[inner sep=0pt,font=\normalsize,text=black] at (156.000,377.000) {3 cm};\end{scope}
\begin{scope}[scale=.6666667,draw=black,text=black]\path[draw=black,line width=.45pt,fill=bookpink] (102,434) rectangle (210,442);\draw[white,line width=.7pt] (104,436)--(208,436);\end{scope}
\begin{scope}[scale=.6666667,draw=black,text=black]\node[inner sep=0pt,font=\normalsize,text=black] at (156.000,421.000) {3 cm};\end{scope}
\begin{scope}[scale=.6666667,draw=black,text=black]\path[draw=black,line width=.45pt,fill=bookpink] (102,478) rectangle (246,486);\draw[white,line width=.7pt] (104,480)--(244,480);\end{scope}
\begin{scope}[scale=.6666667,draw=black,text=black]\node[inner sep=0pt,font=\normalsize,text=black] at (174.000,465.000) {4 cm};\end{scope}
\begin{scope}[scale=.6666667,draw=black,text=black]\path[draw=black,line width=.45pt,fill=bookpink] (344,390) rectangle (416,398);\draw[white,line width=.7pt] (346,392)--(414,392);\end{scope}
\begin{scope}[scale=.6666667,draw=black,text=black]\node[inner sep=0pt,font=\normalsize,text=black] at (380.000,377.000) {2 cm};\end{scope}
\begin{scope}[scale=.6666667,draw=black,text=black]\path[draw=black,line width=.45pt,fill=bookpink] (344,434) rectangle (416,442);\draw[white,line width=.7pt] (346,436)--(414,436);\end{scope}
\begin{scope}[scale=.6666667,draw=black,text=black]\node[inner sep=0pt,font=\normalsize,text=black] at (380.000,421.000) {2 cm};\end{scope}
\begin{scope}[scale=.6666667,draw=black,text=black]\path[draw=black,line width=.45pt,fill=bookpink] (344,478) rectangle (488,486);\draw[white,line width=.7pt] (346,480)--(486,480);\end{scope}
\begin{scope}[scale=.6666667,draw=black,text=black]\node[inner sep=0pt,font=\normalsize,text=black] at (416.000,465.000) {4 cm};\end{scope}
\begin{scope}[scale=.6666667,draw=black,text=black]\path[draw=black,line width=.45pt,fill=bookpink] (590,389) rectangle (662,397);\draw[white,line width=.7pt] (592,391)--(660,391);\end{scope}
\begin{scope}[scale=.6666667,draw=black,text=black]\node[inner sep=0pt,font=\normalsize,text=black] at (626.000,376.000) {2 cm};\end{scope}
\begin{scope}[scale=.6666667,draw=black,text=black]\path[draw=black,line width=.45pt,fill=bookpink] (590,434) rectangle (698,442);\draw[white,line width=.7pt] (592,436)--(696,436);\end{scope}
\begin{scope}[scale=.6666667,draw=black,text=black]\node[inner sep=0pt,font=\normalsize,text=black] at (644.000,421.000) {3 cm};\end{scope}
\begin{scope}[scale=.6666667,draw=black,text=black]\path[draw=black,line width=.45pt,fill=bookpink] (590,478) rectangle (806,486);\draw[white,line width=.7pt] (592,480)--(804,480);\end{scope}
\begin{scope}[scale=.6666667,draw=black,text=black]\node[inner sep=0pt,font=\normalsize,text=black] at (698.000,465.000) {6 cm};\end{scope}
\q{140}{498}{28}{13}{（}
\q{191}{498}{25}{13}{）}
\q{383}{498}{28}{13}{（}
\q{434}{498}{25}{13}{）}
\q{666}{498}{28}{13}{（}
\q{717}{498}{25}{13}{）}
\hdsection{3}{564}{算出下面三角形中$\angle3$的度数。}
\begin{scope}[scale=.6666667,draw=black,text=black]\draw[hdoutline] (134.000,744.000)--(203.500,623.622)--(273.000,744.000)--cycle;\end{scope}
\begin{scope}[scale=.6666667,draw=black,text=black]\draw[line width=.45pt] (195.500,637.478) arc[start angle=119.99990,end angle=60.00010,radius=16];\end{scope}
\begin{scope}[scale=.6666667,draw=black,text=black]\node[inner sep=0pt,font=\normalsize,text=black] at (203.500,654.622) {1};\end{scope}
\begin{scope}[scale=.6666667,draw=black,text=black]\draw[line width=.45pt] (142.000,730.144) arc[start angle=-60.00010,end angle=0.00000,radius=16];\end{scope}
\begin{scope}[scale=.6666667,draw=black,text=black]\node[inner sep=0pt,font=\normalsize,text=black] at (162.579,727.500) {2};\end{scope}
\begin{scope}[scale=.6666667,draw=black,text=black]\draw[line width=.45pt] (257.000,744.000) arc[start angle=180.00000,end angle=240.00010,radius=16];\end{scope}
\begin{scope}[scale=.6666667,draw=black,text=black]\node[inner sep=0pt,font=\normalsize,text=black] at (247.019,729.000) {3};\end{scope}
\begin{scope}[scale=.6666667,draw=black,text=black]\draw[hdoutline] (518.000,727.000)--(604.500,623.910)--(691.000,727.000)--cycle;\end{scope}
\begin{scope}[scale=.6666667,draw=black,text=black]\draw[line width=.45pt] (594.216,636.167) arc[start angle=129.99909,end angle=50.00091,radius=16];\end{scope}
\begin{scope}[scale=.6666667,draw=black,text=black]\node[inner sep=0pt,font=\normalsize,text=black] at (604.500,654.910) {1};\end{scope}
\begin{scope}[scale=.6666667,draw=black,text=black]\draw[line width=.45pt] (528.284,714.743) arc[start angle=-50.00091,end angle=0.00000,radius=16];\end{scope}
\begin{scope}[scale=.6666667,draw=black,text=black]\node[inner sep=0pt,font=\normalsize,text=black] at (547.908,713.053) {2};\end{scope}
\begin{scope}[scale=.6666667,draw=black,text=black]\draw[line width=.45pt] (675.000,727.000) arc[start angle=180.00000,end angle=230.00091,radius=16];\end{scope}
\begin{scope}[scale=.6666667,draw=black,text=black]\node[inner sep=0pt,font=\normalsize,text=black] at (663.811,714.321) {3};\end{scope}
\begin{scope}[scale=.6666667,draw=black,text=black]\draw[hdoutline] (138.000,1017.000)--(138.000,897.000)--(258.000,1017.000)--cycle;\end{scope}
\begin{scope}[scale=.6666667,draw=black,text=black]\draw[line width=.45pt] (138.000,1008.000)--(147.000,1008.000)--(147.000,1017.000);\end{scope}
\begin{scope}[scale=.6666667,draw=black,text=black]\node[inner sep=0pt,font=\normalsize,text=black] at (157.000,999.000) {1};\end{scope}
\begin{scope}[scale=.6666667,draw=black,text=black]\draw[line width=.45pt] (138.000,913.000) arc[start angle=90.00000,end angle=45.00000,radius=16];\end{scope}
\begin{scope}[scale=.6666667,draw=black,text=black]\node[inner sep=0pt,font=\normalsize,text=black] at (151.011,928.412) {2};\end{scope}
\begin{scope}[scale=.6666667,draw=black,text=black]\draw[line width=.45pt] (242.000,1017.000) arc[start angle=180.00000,end angle=225.00000,radius=16];\end{scope}
\begin{scope}[scale=.6666667,draw=black,text=black]\node[inner sep=0pt,font=\normalsize,text=black] at (230.284,1005.519) {3};\end{scope}
\begin{scope}[scale=.6666667,draw=black,text=black]\draw[hdoutline] (511.000,971.000)--(636.312,898.652)--(697.000,971.000)--cycle;\end{scope}
\begin{scope}[scale=.6666667,draw=black,text=black]\draw[line width=.45pt] (622.456,906.652) arc[start angle=150.00031,end angle=50.00891,radius=16];\end{scope}
\begin{scope}[scale=.6666667,draw=black,text=black]\node[inner sep=0pt,font=\normalsize,text=black] at (630.926,929.181) {1};\end{scope}
\begin{scope}[scale=.6666667,draw=black,text=black]\draw[line width=.45pt] (524.856,963.000) arc[start angle=-29.99969,end angle=-0.00000,radius=16];\end{scope}
\begin{scope}[scale=.6666667,draw=black,text=black]\node[inner sep=0pt,font=\normalsize,text=black] at (542.876,962.459) {2};\end{scope}
\begin{scope}[scale=.6666667,draw=black,text=black]\draw[line width=.45pt] (681.000,971.000) arc[start angle=180.00000,end angle=230.00891,radius=16];\end{scope}
\begin{scope}[scale=.6666667,draw=black,text=black]\node[inner sep=0pt,font=\normalsize,text=black] at (669.812,958.319) {3};\end{scope}
\begin{scope}[scale=.6666667,draw=black,text=black]\node[inner sep=0pt,font=\normalsize,text=black] at (112.000,629.000) {①};\end{scope}
\begin{scope}[scale=.6666667,draw=black,text=black]\node[inner sep=0pt,font=\normalsize,text=black] at (489.000,629.000) {②};\end{scope}
\begin{scope}[scale=.6666667,draw=black,text=black]\node[inner sep=0pt,font=\normalsize,text=black] at (111.000,900.000) {③};\end{scope}
\begin{scope}[scale=.6666667,draw=black,text=black]\node[inner sep=0pt,font=\normalsize,text=black] at (489.000,900.000) {④};\end{scope}
\q{135}{784}{325}{13}{$\angle1=60^\circ,\ \angle2=60^\circ,$}
\q{511}{784}{325}{13}{$\angle1=80^\circ,\ \angle2=50^\circ,$}
\q{135}{1057}{325}{13}{$\angle1=90^\circ,\ \angle2=45^\circ,$}
\q{511}{1057}{325}{13}{$\angle1=100^\circ,\ \angle2=30^\circ,$}
\q{135}{839}{75}{13}{$\angle3=$}
\q{178}{839}{28}{13}{（}
\q{250}{839}{25}{13}{）}
\q{265}{839}{25}{12}{。}
\q{511}{839}{75}{13}{$\angle3=$}
\q{554}{839}{28}{13}{（}
\q{626}{839}{25}{13}{）}
\q{641}{839}{25}{12}{。}
\q{135}{1115}{75}{13}{$\angle3=$}
\q{178}{1115}{28}{13}{（}
\q{250}{1115}{25}{13}{）}
\q{265}{1115}{25}{12}{。}
\q{511}{1115}{75}{13}{$\angle3=$}
\q{554}{1115}{28}{13}{（}
\q{626}{1115}{25}{13}{）}
\q{641}{1115}{25}{12}{。}

\begin{scope}[color=answerred,draw=answerred,text=answerred]\ans{254.67}{130.00}{33.33}{10}{$\vphantom{0}\checkmark$}
\ans{505.33}{130.00}{33.33}{10}{$\vphantom{0}\times$}
\ans{254.67}{166.67}{33.33}{10}{$\vphantom{0}\checkmark$}
\ans{505.33}{166.67}{33.33}{10}{$\vphantom{0}\times$}
\ans{107.33}{332.00}{33.33}{10}{$\vphantom{0}\checkmark$}
\ans{270.00}{332.00}{33.33}{10}{$\vphantom{0}\times$}
\ans{458.00}{332.00}{33.33}{10}{$\vphantom{0}\times$}
\ans{139.33}{558.67}{46.67}{11}{$\vphantom{0}60^\circ$}
\ans{388.00}{558.67}{46.67}{11}{$\vphantom{0}50^\circ$}
\ans{138.67}{742.67}{46.67}{11}{$\vphantom{0}45^\circ$}
\ans{388.00}{742.67}{46.67}{11}{$\vphantom{0}50^\circ$}




\ans{188.00}{427.33}{116.00}{9.3}{$\setlength{\jot}{-2pt}\begin{aligned}[t]&\phantom{={}}180^\circ-60^\circ-60^\circ\\&=120^\circ-60^\circ\\&=60^\circ\end{aligned}$}
\ans{467.33}{424.67}{116.00}{9.3}{$\setlength{\jot}{-2pt}\begin{aligned}[t]&\phantom{={}}180^\circ-80^\circ-50^\circ\\&=100^\circ-50^\circ\\&=50^\circ\end{aligned}$}
\ans{188.00}{612.67}{116.00}{9.3}{$\setlength{\jot}{-2pt}\begin{aligned}[t]&\phantom{={}}180^\circ-90^\circ-45^\circ\\&=90^\circ-45^\circ\\&=45^\circ\end{aligned}$}
\ans{467.33}{616.67}{116.00}{9.3}{$\setlength{\jot}{-2pt}\begin{aligned}[t]&\phantom{={}}180^\circ-100^\circ-30^\circ\\&=80^\circ-30^\circ\\&=50^\circ\end{aligned}$}
\end{scope}
\end{scope}\pgfresetboundingbox\path[use as bounding box] (56.0,242.0) rectangle (543.0,374.0);
\end{tikzpicture}}\par



\par





\par\noindent

\QHeading{第1、2题　组成三角形的另一种判断方法}

\Answer{除“较短两边的和大于最长边”外，也可以用“最长边减去次长边，所得差小于最短边”来判断。例如3、4、6：6减4得2，小于3，可以组成；1、2、3：3减2得1，等于1，不可以组成。两种判断是同一个三角形边长条件的不同使用方式。}

\QHeading{第3题　算出下面三角形中$\angle3$的度数。}

\par\noindent\makebox[\linewidth][c]{
\begin{tikzpicture}[x=1bp,y=-1bp,color=black,draw=black,text=black]
\begin{scope}\clip (62.0,406.0) rectangle (557.0,762.0);
\hdhead{42}{多边形的认识——信息窗2}
\hdsection{1}{140}{能围成三角形的在（　）里画“$\checkmark$”。}
\q{103}{197}{270}{12.5}{① 3 cm　4 cm　6 cm}
\q{480}{197}{270}{12.5}{② 1 cm　2 cm　3 cm}
\q{103}{251}{270}{12.5}{③ 5 cm　7 cm　11 cm}
\q{480}{251}{270}{12.5}{④ 12 m　15 m　29 m}
\q{355}{195}{28}{13}{（}
\q{409}{195}{25}{13}{）}
\q{736}{195}{28}{13}{（}
\q{790}{195}{25}{13}{）}
\q{355}{250}{28}{13}{（}
\q{409}{250}{25}{13}{）}
\q{736}{250}{28}{13}{（}
\q{790}{250}{25}{13}{）}
\hdsection{2}{314}{下面每组中的3根小棒能围成三角形吗? 能的在（　）里画“$\checkmark$”。}
\begin{scope}[scale=.6666667,draw=black,text=black]\path[draw=black,line width=.45pt,fill=bookpink] (102,390) rectangle (210,398);\draw[white,line width=.7pt] (104,392)--(208,392);\end{scope}
\begin{scope}[scale=.6666667,draw=black,text=black]\node[inner sep=0pt,font=\normalsize,text=black] at (156.000,377.000) {3 cm};\end{scope}
\begin{scope}[scale=.6666667,draw=black,text=black]\path[draw=black,line width=.45pt,fill=bookpink] (102,434) rectangle (210,442);\draw[white,line width=.7pt] (104,436)--(208,436);\end{scope}
\begin{scope}[scale=.6666667,draw=black,text=black]\node[inner sep=0pt,font=\normalsize,text=black] at (156.000,421.000) {3 cm};\end{scope}
\begin{scope}[scale=.6666667,draw=black,text=black]\path[draw=black,line width=.45pt,fill=bookpink] (102,478) rectangle (246,486);\draw[white,line width=.7pt] (104,480)--(244,480);\end{scope}
\begin{scope}[scale=.6666667,draw=black,text=black]\node[inner sep=0pt,font=\normalsize,text=black] at (174.000,465.000) {4 cm};\end{scope}
\begin{scope}[scale=.6666667,draw=black,text=black]\path[draw=black,line width=.45pt,fill=bookpink] (344,390) rectangle (416,398);\draw[white,line width=.7pt] (346,392)--(414,392);\end{scope}
\begin{scope}[scale=.6666667,draw=black,text=black]\node[inner sep=0pt,font=\normalsize,text=black] at (380.000,377.000) {2 cm};\end{scope}
\begin{scope}[scale=.6666667,draw=black,text=black]\path[draw=black,line width=.45pt,fill=bookpink] (344,434) rectangle (416,442);\draw[white,line width=.7pt] (346,436)--(414,436);\end{scope}
\begin{scope}[scale=.6666667,draw=black,text=black]\node[inner sep=0pt,font=\normalsize,text=black] at (380.000,421.000) {2 cm};\end{scope}
\begin{scope}[scale=.6666667,draw=black,text=black]\path[draw=black,line width=.45pt,fill=bookpink] (344,478) rectangle (488,486);\draw[white,line width=.7pt] (346,480)--(486,480);\end{scope}
\begin{scope}[scale=.6666667,draw=black,text=black]\node[inner sep=0pt,font=\normalsize,text=black] at (416.000,465.000) {4 cm};\end{scope}
\begin{scope}[scale=.6666667,draw=black,text=black]\path[draw=black,line width=.45pt,fill=bookpink] (590,389) rectangle (662,397);\draw[white,line width=.7pt] (592,391)--(660,391);\end{scope}
\begin{scope}[scale=.6666667,draw=black,text=black]\node[inner sep=0pt,font=\normalsize,text=black] at (626.000,376.000) {2 cm};\end{scope}
\begin{scope}[scale=.6666667,draw=black,text=black]\path[draw=black,line width=.45pt,fill=bookpink] (590,434) rectangle (698,442);\draw[white,line width=.7pt] (592,436)--(696,436);\end{scope}
\begin{scope}[scale=.6666667,draw=black,text=black]\node[inner sep=0pt,font=\normalsize,text=black] at (644.000,421.000) {3 cm};\end{scope}
\begin{scope}[scale=.6666667,draw=black,text=black]\path[draw=black,line width=.45pt,fill=bookpink] (590,478) rectangle (806,486);\draw[white,line width=.7pt] (592,480)--(804,480);\end{scope}
\begin{scope}[scale=.6666667,draw=black,text=black]\node[inner sep=0pt,font=\normalsize,text=black] at (698.000,465.000) {6 cm};\end{scope}
\q{140}{498}{28}{13}{（}
\q{191}{498}{25}{13}{）}
\q{383}{498}{28}{13}{（}
\q{434}{498}{25}{13}{）}
\q{666}{498}{28}{13}{（}
\q{717}{498}{25}{13}{）}
\hdsection{3}{564}{算出下面三角形中$\angle3$的度数。}
\begin{scope}[scale=.6666667,draw=black,text=black]\draw[hdoutline] (134.000,744.000)--(203.500,623.622)--(273.000,744.000)--cycle;\end{scope}
\begin{scope}[scale=.6666667,draw=black,text=black]\draw[line width=.45pt] (195.500,637.478) arc[start angle=119.99990,end angle=60.00010,radius=16];\end{scope}
\begin{scope}[scale=.6666667,draw=black,text=black]\node[inner sep=0pt,font=\normalsize,text=black] at (203.500,654.622) {1};\end{scope}
\begin{scope}[scale=.6666667,draw=black,text=black]\draw[line width=.45pt] (142.000,730.144) arc[start angle=-60.00010,end angle=0.00000,radius=16];\end{scope}
\begin{scope}[scale=.6666667,draw=black,text=black]\node[inner sep=0pt,font=\normalsize,text=black] at (162.579,727.500) {2};\end{scope}
\begin{scope}[scale=.6666667,draw=black,text=black]\draw[line width=.45pt] (257.000,744.000) arc[start angle=180.00000,end angle=240.00010,radius=16];\end{scope}
\begin{scope}[scale=.6666667,draw=black,text=black]\node[inner sep=0pt,font=\normalsize,text=black] at (247.019,729.000) {3};\end{scope}
\begin{scope}[scale=.6666667,draw=black,text=black]\draw[hdoutline] (518.000,727.000)--(604.500,623.910)--(691.000,727.000)--cycle;\end{scope}
\begin{scope}[scale=.6666667,draw=black,text=black]\draw[line width=.45pt] (594.216,636.167) arc[start angle=129.99909,end angle=50.00091,radius=16];\end{scope}
\begin{scope}[scale=.6666667,draw=black,text=black]\node[inner sep=0pt,font=\normalsize,text=black] at (604.500,654.910) {1};\end{scope}
\begin{scope}[scale=.6666667,draw=black,text=black]\draw[line width=.45pt] (528.284,714.743) arc[start angle=-50.00091,end angle=0.00000,radius=16];\end{scope}
\begin{scope}[scale=.6666667,draw=black,text=black]\node[inner sep=0pt,font=\normalsize,text=black] at (547.908,713.053) {2};\end{scope}
\begin{scope}[scale=.6666667,draw=black,text=black]\draw[line width=.45pt] (675.000,727.000) arc[start angle=180.00000,end angle=230.00091,radius=16];\end{scope}
\begin{scope}[scale=.6666667,draw=black,text=black]\node[inner sep=0pt,font=\normalsize,text=black] at (663.811,714.321) {3};\end{scope}
\begin{scope}[scale=.6666667,draw=black,text=black]\draw[hdoutline] (138.000,1017.000)--(138.000,897.000)--(258.000,1017.000)--cycle;\end{scope}
\begin{scope}[scale=.6666667,draw=black,text=black]\draw[line width=.45pt] (138.000,1008.000)--(147.000,1008.000)--(147.000,1017.000);\end{scope}
\begin{scope}[scale=.6666667,draw=black,text=black]\node[inner sep=0pt,font=\normalsize,text=black] at (157.000,999.000) {1};\end{scope}
\begin{scope}[scale=.6666667,draw=black,text=black]\draw[line width=.45pt] (138.000,913.000) arc[start angle=90.00000,end angle=45.00000,radius=16];\end{scope}
\begin{scope}[scale=.6666667,draw=black,text=black]\node[inner sep=0pt,font=\normalsize,text=black] at (151.011,928.412) {2};\end{scope}
\begin{scope}[scale=.6666667,draw=black,text=black]\draw[line width=.45pt] (242.000,1017.000) arc[start angle=180.00000,end angle=225.00000,radius=16];\end{scope}
\begin{scope}[scale=.6666667,draw=black,text=black]\node[inner sep=0pt,font=\normalsize,text=black] at (230.284,1005.519) {3};\end{scope}
\begin{scope}[scale=.6666667,draw=black,text=black]\draw[hdoutline] (511.000,971.000)--(636.312,898.652)--(697.000,971.000)--cycle;\end{scope}
\begin{scope}[scale=.6666667,draw=black,text=black]\draw[line width=.45pt] (622.456,906.652) arc[start angle=150.00031,end angle=50.00891,radius=16];\end{scope}
\begin{scope}[scale=.6666667,draw=black,text=black]\node[inner sep=0pt,font=\normalsize,text=black] at (630.926,929.181) {1};\end{scope}
\begin{scope}[scale=.6666667,draw=black,text=black]\draw[line width=.45pt] (524.856,963.000) arc[start angle=-29.99969,end angle=-0.00000,radius=16];\end{scope}
\begin{scope}[scale=.6666667,draw=black,text=black]\node[inner sep=0pt,font=\normalsize,text=black] at (542.876,962.459) {2};\end{scope}
\begin{scope}[scale=.6666667,draw=black,text=black]\draw[line width=.45pt] (681.000,971.000) arc[start angle=180.00000,end angle=230.00891,radius=16];\end{scope}
\begin{scope}[scale=.6666667,draw=black,text=black]\node[inner sep=0pt,font=\normalsize,text=black] at (669.812,958.319) {3};\end{scope}
\begin{scope}[scale=.6666667,draw=black,text=black]\node[inner sep=0pt,font=\normalsize,text=black] at (112.000,629.000) {①};\end{scope}
\begin{scope}[scale=.6666667,draw=black,text=black]\node[inner sep=0pt,font=\normalsize,text=black] at (489.000,629.000) {②};\end{scope}
\begin{scope}[scale=.6666667,draw=black,text=black]\node[inner sep=0pt,font=\normalsize,text=black] at (111.000,900.000) {③};\end{scope}
\begin{scope}[scale=.6666667,draw=black,text=black]\node[inner sep=0pt,font=\normalsize,text=black] at (489.000,900.000) {④};\end{scope}
\q{135}{784}{325}{13}{$\angle1=60^\circ,\ \angle2=60^\circ,$}
\q{511}{784}{325}{13}{$\angle1=80^\circ,\ \angle2=50^\circ,$}
\q{135}{1057}{325}{13}{$\angle1=90^\circ,\ \angle2=45^\circ,$}
\q{511}{1057}{325}{13}{$\angle1=100^\circ,\ \angle2=30^\circ,$}
\q{135}{839}{75}{13}{$\angle3=$}
\q{178}{839}{28}{13}{（}
\q{250}{839}{25}{13}{）}
\q{265}{839}{25}{12}{。}
\q{511}{839}{75}{13}{$\angle3=$}
\q{554}{839}{28}{13}{（}
\q{626}{839}{25}{13}{）}
\q{641}{839}{25}{12}{。}
\q{135}{1115}{75}{13}{$\angle3=$}
\q{178}{1115}{28}{13}{（}
\q{250}{1115}{25}{13}{）}
\q{265}{1115}{25}{12}{。}
\q{511}{1115}{75}{13}{$\angle3=$}
\q{554}{1115}{28}{13}{（}
\q{626}{1115}{25}{13}{）}
\q{641}{1115}{25}{12}{。}

\begin{scope}[color=answerred,draw=answerred,text=answerred]\ans{254.67}{130.00}{33.33}{10}{$\vphantom{0}\checkmark$}
\ans{505.33}{130.00}{33.33}{10}{$\vphantom{0}\times$}
\ans{254.67}{166.67}{33.33}{10}{$\vphantom{0}\checkmark$}
\ans{505.33}{166.67}{33.33}{10}{$\vphantom{0}\times$}
\ans{107.33}{332.00}{33.33}{10}{$\vphantom{0}\checkmark$}
\ans{270.00}{332.00}{33.33}{10}{$\vphantom{0}\times$}
\ans{458.00}{332.00}{33.33}{10}{$\vphantom{0}\times$}
\ans{139.33}{558.67}{46.67}{11}{$\vphantom{0}60^\circ$}
\ans{388.00}{558.67}{46.67}{11}{$\vphantom{0}50^\circ$}
\ans{138.67}{742.67}{46.67}{11}{$\vphantom{0}45^\circ$}
\ans{388.00}{742.67}{46.67}{11}{$\vphantom{0}50^\circ$}




\ans{188.00}{427.33}{116.00}{9.3}{$\setlength{\jot}{-2pt}\begin{aligned}[t]&\phantom{={}}180^\circ-60^\circ-60^\circ\\&=120^\circ-60^\circ\\&=60^\circ\end{aligned}$}
\ans{467.33}{424.67}{116.00}{9.3}{$\setlength{\jot}{-2pt}\begin{aligned}[t]&\phantom{={}}180^\circ-80^\circ-50^\circ\\&=100^\circ-50^\circ\\&=50^\circ\end{aligned}$}
\ans{188.00}{612.67}{116.00}{9.3}{$\setlength{\jot}{-2pt}\begin{aligned}[t]&\phantom{={}}180^\circ-90^\circ-45^\circ\\&=90^\circ-45^\circ\\&=45^\circ\end{aligned}$}
\ans{467.33}{616.67}{116.00}{9.3}{$\setlength{\jot}{-2pt}\begin{aligned}[t]&\phantom{={}}180^\circ-100^\circ-30^\circ\\&=80^\circ-30^\circ\\&=50^\circ\end{aligned}$}
\end{scope}
\end{scope}\pgfresetboundingbox\path[use as bounding box] (62.0,406.0) rectangle (557.0,762.0);
\end{tikzpicture}}\par



\par









\QHeading{第3题（1）　另一种列式}

\Answer{先把两个已知角相加，再从180度里减去。
\Formula{\begin{aligned}[t]&\phantom{={}}180-(60+60)\\
&=180-120\\
&=60\text{（度）}\end{aligned}}\par 答：角3是60度。}

\QHeading{第3题（2）　另一种列式}

\Answer{先把两个已知角相加，再从180度里减去。
\Formula{\begin{aligned}[t]&\phantom{={}}180-(80+50)\\
&=180-130\\
&=50\text{（度）}\end{aligned}}\par 答：角3是50度。}

\QHeading{第3题（3）　另一种列式}

\Answer{先把两个已知角相加，再从180度里减去。
\Formula{\begin{aligned}[t]&\phantom{={}}180-(90+45)\\
&=180-135\\
&=45\text{（度）}\end{aligned}}\par 答：角3是45度。}

\QHeading{第3题（4）　另一种列式}

\Answer{先把两个已知角相加，再从180度里减去。
\Formula{\begin{aligned}[t]&\phantom{={}}180-(100+30)\\
&=180-130\\
&=50\text{（度）}\end{aligned}}\par 答：角3是50度。}
